\documentclass[11pt]{article}
\usepackage[margin=0.8in]{geometry}
\usepackage[T1]{fontenc}
\usepackage{lmodern,amsmath,amssymb,amsthm,booktabs,array}
\usepackage[colorlinks=true,linkcolor=blue,urlcolor=blue,citecolor=blue]{hyperref}
\hypersetup{pdftitle={Eight sequence assertions: negabinary, divisor sums, and tetration},pdfauthor={Blake Stone},pdfsubject={Seven proofs and one disproof with complete Lean verification}}
\setlength{\parindent}{0pt}
\setlength{\parskip}{0.4em}
\newtheorem{theorem}{Theorem}[section]
\newtheorem{lemma}[theorem]{Lemma}
\newtheorem{corollary}[theorem]{Corollary}
\newcommand{\N}{\mathbb N}
\newcommand{\Z}{\mathbb Z}
\begin{document}
\begin{center}
{\LARGE\bfseries Eight sequence assertions:\\[3pt]
negabinary, divisor sums, and tetration}\par
\vspace{0.5em}
{\large Blake Stone}\par
Proof developed with OpenAI Astra 6.2\par
November 26, 2026
\end{center}

\textbf{Abstract.}
We resolve eight assertions recorded as conjectures in five OEIS entries.
Seven are proved and one is refuted. The results comprise a negabinary
exclusive-OR identity, a restriction on integers with two squarefree
divisor sums, two tetration congruences and a counterexample to a third,
two divisor-count identities for a recursive triangle, and existence of
palindromic proper-divisor residue lists at every positive index.
All eight assertions are checked in Lean from their original definitions;
the six accompanying proof modules have fresh package builds and
independent compilations. The arguments use classical ingredients,
including Kehowski's squarefree-divisor obstruction, modular exponential
lifting, Uchimura's divisor identity, and primes congruent to one.
Those ingredients are credited. The contribution is the explicit
resolution and formal verification of the selected statements, with
bounded prior-work checks and no claim of exhaustive historical priority.

\begin{center}
\begin{tabular}{@{}llp{3.7in}@{}}
\toprule
Entry & Assertions & Result\\
\midrule
A016789 & 1 proved & The exact three-way negabinary XOR identity.\\
A300984 & 1 proved & The prime-power shape, stronger necessary conditions, and a lower bound.\\
A324017 & 2 proved, 1 false & Fixed points and quotient residues; a canonical counterexample to the next-digit claim.\\
A389524 & 2 proved & Row sums $d(n)$ and diagonal $n-d(n)$.\\
A393833 & 1 proved & Existence, with arbitrarily large prime-power witnesses.\\
\bottomrule
\end{tabular}
\end{center}

\section{A016789: a negabinary exclusive-OR identity}

For $m\ge0$, let $P(m)$ be the integer obtained by writing $m$ in base
$-2$ and reading the same digits in base $2$. Define $N(m)$ similarly
for $-m$. These are A005351 and A005352, with the latter extended by
$N(0)=0$. The standard representation and conversion algorithms are
described by Arndt~\cite{Arndt}. Gilian Breysens's July 21, 2017
conjecture in A016789~\cite{A016789} is the following assertion, where
$\oplus$ denotes bitwise exclusive OR.

\begin{theorem}\label{thm:xor}
For every $n\ge0$,
\[
 n\oplus P(n+1)\oplus N(n+1)=3n+2.
\]
\end{theorem}

\begin{proof}
For a binary code $c=\sum_i b_i2^i$, with $b_i\in\{0,1\}$, put
$V(c)=\sum_i b_i(-2)^i$. Appending a low bit gives
\begin{equation}\label{eq:value}
 V(2c+b)=b-2V(c),\qquad b\in\{0,1\}.
\end{equation}
This map is injective. Equal values force equal low bits by reduction
modulo two. Removing those bits gives equal values for the two truncated
codes. Repeated division by two terminates, so the original codes agree.

Existence of codes for $m$ and $-m$ follows simultaneously by strong
induction. For $m=0,1$ use the codes $0,1,3$. For $m\ge2$, set
$b=m\bmod2$, $u=\lfloor m/2\rfloor$, and $v=\lceil m/2\rceil$.
Both $u,v<m$. A code $c$ for $-u$ gives the code $2c+b$ for $m$, and
a code $d$ for $v$ gives the code $2d+b$ for $-m$, by~\eqref{eq:value}.
Thus $P$ and $N$ are the unique codes of their stated signed values.

Evaluation and uniqueness now give all four recurrences
\begin{equation}\label{eq:code-recurrences}
\begin{aligned}
P(2t)&=2N(t), &N(2t)&=2P(t),\\
P(2t+1)&=2N(t)+1, &N(2t+1)&=2P(t+1)+1.
\end{aligned}
\end{equation}
Define three expressions
\[
\begin{aligned}
S_0(t)&=t\oplus P(t)\oplus N(t),\\
S_1(t)&=t\oplus N(t)\oplus P(t+1),\\
S_2(t)&=t\oplus P(t+1)\oplus N(t+1).
\end{aligned}
\]
For low bits $b,c\in\{0,1\}$,
\[
(2x+b)\oplus(2y+c)=2(x\oplus y)+(b\oplus c).
\]
Combining this with~\eqref{eq:code-recurrences} gives
\begin{center}
\begin{tabular}{@{}lcc@{}}
\toprule
 & At $2t$ & At $2t+1$\\
\midrule
$S_0$ & $2S_0(t)$ & $2S_1(t)+1$\\
$S_1$ & $2S_0(t)+1$ & $2S_2(t)$\\
$S_2$ & $2S_1(t)$ & $2S_2(t)+1$\\
\bottomrule
\end{tabular}
\end{center}
For example,
\[
S_2(2t)=(2t)\oplus(2N(t)+1)\oplus(2P(t+1)+1)=2S_1(t).
\]
The initial values $P(0)=N(0)=0$, $P(1)=1$, $N(1)=3$ give
$S_0(0)=0$, $S_1(0)=1$, $S_2(0)=2$. The table proves
$S_i(t)=3t+i$ for $i=0,1,2$ simultaneously by strong induction:
at a positive index, every branch uses only its strictly smaller half.
The $i=2$ case is the required identity.
\end{proof}

The two companion identities are parts of this proof, not additional
recorded conjectures. In Lean, $V$ is literally
\texttt{Nat.ofDigits (-2 : $\Z$) (Nat.digits 2 code)}. Existence and
uniqueness are proved before the code functions are defined.
\texttt{original\_conjecture} states the result directly for any two
codes whose base-$-2$ values are $n+1$ and $-(n+1)$.

\section{A300984: two squarefree divisor sums}

For a positive integer $n$, define
\[
 S(n)=\sum_{\substack{d\mid n\\d\text{ squarefree}}}d,
 \qquad
 T(n)=\sum_{\substack{d\mid n\\d\text{ not squarefree}}}d.
\]
Michel Lagneau's March 17, 2018 conjecture in A300984~\cite{A300984}
asserts that if both sums are squarefree, then $n=2^i p^j$ for a prime
$p$. We obtain the following stronger necessary conditions.

\begin{theorem}\label{thm:squarefree}
If $S(n)$ and $T(n)$ are squarefree, then
\[
n=2^i p^j,\qquad i,j\ge2,\qquad p\text{ prime},\qquad
p\equiv1\pmod{12},
\]
and $3(p+1)$ is squarefree. In particular, $n\ge676$.
\end{theorem}

\begin{proof}
Let $\mathcal P(n)$ denote the set of distinct prime divisors of $n$.
The classical subset expansion of the squarefree divisors gives
\begin{equation}\label{eq:squarefree-product}
S(n)=\prod_{q\in\mathcal P(n)}(1+q)
     =\sigma(\operatorname{rad}(n)).
\end{equation}
If two distinct odd primes $p,q$ divide $n$, their two factors in this
product are even. Therefore $4\mid S(n)$, contradicting squarefreeness.
There is at most one odd prime divisor of $n$.

Next, squarefreeness of $T(n)$ forces at least two distinct prime squares
to divide $n$. If no prime square divides $n$, all its divisors are
squarefree and $T(n)=0$, which is not squarefree. If exactly one prime
square $r^2$ divides $n$, every nonsquarefree divisor $d$ is divisible by
some prime square $s^2$. Since $s^2\mid d\mid n$, uniqueness forces
$s=r$. Hence $r^2$ divides every summand of $T(n)$ and therefore its
sum, again contradicting squarefreeness.

There must consequently be exactly two distinct prime divisors, namely
$2$ and one odd prime $p$, and both $4$ and $p^2$ divide $n$. Unique
factorization gives $n=2^i p^j$ with $i,j\ge2$.
Equation~\eqref{eq:squarefree-product} now reads $S(n)=3(p+1)$.
Squarefreeness excludes $4\mid p+1$. Since $p$ is odd, this forces
$p\equiv1\pmod4$. It also excludes $3\mid p+1$; otherwise
$9\mid3(p+1)$. The prime $p$ is not $3$, because $p+1$ would then
equal $4$. Thus $p\equiv1\pmod3$ as well. Combining the congruences
gives $p\equiv1\pmod{12}$, whence $p\ge13$ and
$n\ge2^2 13^2=676$.
\end{proof}

These conditions are not asserted to be sufficient: $T(n)$ still has
to be squarefree. The separate conjecture about the frequency of certain
associated primes is outside the result. At $676$, direct calculation
gives $S(676)=42=2\cdot3\cdot7$ and
$T(676)=1239=3\cdot7\cdot59$, consistent with the lower bound.
The formal theorem proves the universal lower bound; the displayed
calculation is a transparent numerical illustration.

\textbf{Prior argument.}
The product formula is recorded by Vladeta Jovovic in A048250 in
2001~\cite{A048250}. Walter Kehowski's 2007 comment on
A065302~\cite{A065302} already proves the corresponding $2p$ shape
and $p\equiv1\pmod{12}$ restriction for squarefree nonprime integers
with squarefree ordinary divisor sum. His two-odd-prime obstruction
applies here to $\operatorname{rad}(n)$ and is explicitly credited.
The second sum supplies the additional exponent restrictions.
Lean begins with the two actual filtered sums over \texttt{Nat.divisors},
proves~\eqref{eq:squarefree-product}, and then proves all the stated
necessary conditions. No factorization characterization is assumed.

\section{A324017: tetration congruences and a false strengthening}

Write $T_0(a)=1$ and $T_{m+1}(a)=a^{T_m(a)}$. This section uses $T$ for
a power tower, independently of the divisor sum in the preceding section.
Davis Smith's A324017~\cite{A324017}, introduced March 28, 2019, is
\[
A(m,n)=T_m(2n-1)\bmod(2n)^m.
\]
Its three conjectural comments concern a fixed-point property, the next
base-$2n$ digit, and a stronger compatibility property.

\begin{theorem}\label{thm:tetration}
Let $b=2n\ge2$ and $a=b-1$.
\begin{enumerate}
\item For $m\ge1$ and $k\ge0$, if $T_m(a)\equiv k\pmod{b^m}$,
then $a^k\equiv k\pmod{b^m}$.
\item For $m\ge2$,
\[
T_m(a)\bmod b^2=b^2-b-1.
\]
Consequently, for the source's domain $n>1,m\ge2$,
\[
\left\lfloor\frac{T_m(2n-1)}{2n}\right\rfloor
\equiv2(n-1)\pmod{2n}.
\]
\item The claim that $T_m(2n-1)\equiv j\pmod{(2n)^{m+1}}$
implies $A(m+1,n)=j$ is false, even with $j$ the least nonnegative
residue. A valid replacement is
\[
A(m+1,n)\bmod(2n)^m=A(m,n),\qquad m,n\ge1.
\]
\end{enumerate}
\end{theorem}

\begin{proof}
We first establish a period that gains one power of the modulus.
For $r\ge1$, if $x\equiv1\pmod{b^r}$, then
\[
x^b-1=(x-1)(1+x+\cdots+x^{b-1})
\]
is divisible by $b^{r+1}$: the first factor is divisible by $b^r$,
and the second is congruent to $b$ modulo $b$.
Since $b$ is even, the binomial theorem gives
\[
(b-1)^b=(-1+b)^b\equiv1-b^2\equiv1\pmod{b^2}.
\]
Repeated application of the preceding lifting argument yields
\begin{equation}\label{eq:period}
a^{b^r}\equiv1\pmod{b^{r+1}}\qquad(r\ge1).
\end{equation}
It follows that
\begin{equation}\label{eq:exp-lift}
u\equiv v\pmod{b^r}\quad\Longrightarrow\quad
a^u\equiv a^v\pmod{b^{r+1}}
\end{equation}
for nonnegative $u,v$: express both as their common remainder plus
a multiple of $b^r$ and use~\eqref{eq:period}.

Every tower $T_m(a)$ is odd. Since $a\equiv-1\pmod b$, every
positive-height tower is congruent to $-1$ modulo $b$.
In particular $T_2(a)\equiv T_1(a)\pmod b$.
Induction using~\eqref{eq:exp-lift} now proves
\begin{equation}\label{eq:stable}
T_{m+1}(a)\equiv T_m(a)\pmod{b^m}\qquad(m\ge1).
\end{equation}
Indeed, exponentiating the statement at height $m$ proves the statement
at height $m+1$. If $T_m(a)\equiv k\pmod{b^m}$, then
\[
a^k\equiv a^{T_m(a)}=T_{m+1}(a)\equiv T_m(a)\equiv k\pmod{b^m},
\]
using~\eqref{eq:exp-lift} followed by~\eqref{eq:stable}. This proves
part 1, including the original ranges $m,n,k>1$.

For part 2, if $E$ is odd, a first-order binomial expansion gives
\[
(b-1)^E\equiv-1+bE\pmod{b^2}.
\]
When $m\ge2$, the exponent $E=T_{m-1}(a)$ is odd and is $-1$ modulo
$b$. Therefore $T_m(a)\equiv-b-1\pmod{b^2}$. Its least nonnegative
residue is
\[
b^2-b-1=b(b-2)+(b-1),
\]
which lies between zero and $b^2-1$, including $b=2$.
For nonnegative $t$,
\[
\left\lfloor t/b\right\rfloor\bmod b
=\left\lfloor (t\bmod b^2)/b\right\rfloor.
\]
Substitution gives $b-2=2(n-1)$ as the quotient residue.

For part 3, choose $n=m=2$ and $j=27$. Then
\[
T_2(3)=3^3=27\equiv27\pmod{64},\qquad
A(3,2)=3^{27}\bmod64=59\ne27.
\]
The residue $27$ is canonical, so the counterexample persists even
after restricting $j$ to $0\le j<64$.
Finally, since $b^m\mid b^{m+1}$, reducing $T_{m+1}(a)$ first modulo
$b^{m+1}$ and then modulo $b^m$ is equivalent to direct reduction
modulo $b^m$. Equation~\eqref{eq:stable} proves the replacement relation.
\end{proof}

The source's examples in the third comment also transpose the array
indices: the displayed table places $563$ at $A(3,6)$ and $16115$ at
$A(4,6)$. The disproof addresses the quantified claim as written and
does not infer a different intended conjecture from those examples.

\textbf{Prior theory and exact formal scope.}
Modular tetration and exponential fixed points are established topics.
The OEIS reference to Trigg's Problem 559 in 1981 leads to a decimal
lifting problem; Shapiro's 1987 solution proves the relevant decimal
lifting and uniqueness~\cite{Crux1981,Crux1987}. General stabilization
and fixed points are treated by Shapiro and Shapiro~\cite{Shapiro},
and a broader lifting framework appears in Jim\'enez Urroz and
Yebra~\cite{Yebra}. The lifting method is credited as classical.
Here it is specialized to every even $b$ and $a=b-1$, with the precise
height bound and the exact statements above fully checked.
The Lean definitions are literal integer tetration and modular reduction;
the proofs do not substitute a remainder-iteration algorithm for the
original towers. The counterexample uses kernel-checked arithmetic.

\section{A389524: row sums and a divisor-count diagonal}

Mats Granvik's A389524~\cite{A389524}, introduced October 7, 2025,
defines an integer triangle. Its source program specifies
\[
T(n,1)=\begin{cases}1,&n=1,\\0,&n\ne1,\end{cases}
\]
and, for $k\ge2$, sets $T(n,k)=0$ when $n<k$, while for $n\ge k$
\begin{equation}\label{eq:triangle}
T(n,k)=1+\sum_{i=1}^{k-1}T(n-i,k-1)
          -\sum_{i=1}^{n-1}T(n-i,k).
\end{equation}
The notation $T$ in this section refers to this triangle. Let $d(n)$
be the number of positive divisors of $n$.

\begin{theorem}\label{thm:triangle}
For every positive integer $n$,
\[
\sum_{k=1}^{n}T(n,k)=d(n),\qquad
T(2n,n)=n-d(n).
\]
\end{theorem}

These are Conjectures 1 and 3 in the source. The separately recorded
weighted orphan-count conjecture, numbered 2, is not part of this result.

\begin{proof}
We derive the column polynomials from the defining recurrence.
Let $C_k(x)=\sum_{n\ge0}T(n,k)x^n$, with the triangle extended by zero
outside its natural support. The first column gives $C_1=x$.
Moving the last sum in~\eqref{eq:triangle} to the left and taking formal
generating series gives, for $k\ge2$,
\[
\frac{C_k(x)}{1-x}
=\frac{x^k}{1-x}+(x+\cdots+x^{k-1})C_{k-1}(x).
\]
Therefore $C_k=x^k+x(1-x^{k-1})C_{k-1}$.
Define polynomials $R_0=0$ and
\begin{equation}\label{eq:R}
R_{k+1}=1+(1-x^k)R_k.
\end{equation}
Then $R_1=1$ and $C_k=x^kR_k$. The Lean proof establishes this bridge
by strong induction on the original row index, using finite coefficient
sums; the generating-series computation is not an assumption.

Put
\[
S_N(x)=\sum_{k=1}^{N}x^kR_k(x),\qquad
W_N(x)=\sum_{k=1}^{N}k x^k\prod_{j=k+1}^{N}(1-x^j).
\]
Equation~\eqref{eq:R} gives $R_{k+1}-R_k=1-x^kR_k$, hence
$R_{N+1}=N+1-S_N$ by telescoping. It follows that
\[
S_{N+1}=(1-x^{N+1})S_N+(N+1)x^{N+1}.
\]
The product expression $W_N$ satisfies the same recurrence and the same
zero initial value. Thus $S_N=W_N$.

We next prove the needed finite coefficient version of Uchimura's
classical divisor identity~\cite{Uchimura,Bressoud}:
\begin{equation}\label{eq:finite-uchimura}
W_N(x)\equiv\sum_{i=1}^{N}\sum_{j=1}^{N}x^{ij}\pmod{x^{N+1}}.
\end{equation}
For a finite algebraic proof, work in
$A=\Z[x]/(x^{N+1})$ and write $q$ for the class of $x$.
In the second polynomial variable $z$, define
\[
H(z)=\sum_{k=0}^{N}z^kq^k\prod_{j=k+1}^{N}(1-q^j),
\quad
G(z)=\prod_{j=1}^{N}(1-zq^j),
\quad
P=\prod_{j=1}^{N}(1-q^j).
\]
Finite coefficient comparison and cancellation of the product endpoints
give
\[
H(qz)=(1-qz)H(z),\qquad (1-qz)G(qz)=G(z).
\]
The boundary terms vanish because $q^{N+1}=0$.
As $1-qz$ is a unit in $A[z]$, their product $K(z)=H(z)G(z)$ satisfies
$K(qz)=K(z)$. For every $r>0$, comparison of the coefficient of $z^r$
gives $(1-q^r)[z^r]K=0$. The factor $1-q^r$ is a unit, since $q^r$
is nilpotent. Thus $K$ is constant, and its constant coefficient is $P$:
\begin{equation}\label{eq:HG}
H(z)G(z)=P.
\end{equation}
At $z=1$, $G(1)=P$ is a unit, so $H(1)=1$. Differentiate
\eqref{eq:HG} with respect to $z$ and evaluate at one. The product rule
gives
\[
G'(1)=-P\sum_{i=1}^{N}\frac{q^i}{1-q^i},
\]
and hence
\[
W_N(q)=H'(1)=\sum_{i=1}^{N}\frac{q^i}{1-q^i}
             =\sum_{i=1}^{N}\sum_{j=1}^{N}q^{ij}.
\]
The fractions mean multiplication by units' inverses. Their geometric
expansions are finite because all terms of degree at least $N+1$ vanish.
This proves~\eqref{eq:finite-uchimura} without an analytic limit.

Take $N=n\ge1$. The coefficient of $x^n$ on the right counts the pairs
$1\le i,j\le n$ with $ij=n$. Every divisor $i$ of $n$ gives exactly
one such pair $(i,n/i)$, so the coefficient is $d(n)$. Since $S_n=W_n$
and $C_k=x^kR_k$,
\[
\sum_{k=1}^{n}T(n,k)=[x^n]S_n=d(n).
\]
All positive columns beyond $n$ are zero, so this is also the source's
sum over all positive columns.

Finally, $R_n(0)=n$ follows from~\eqref{eq:R}. The identities
\[
R_n=n-S_{n-1},\qquad S_n=S_{n-1}+x^nR_n
\]
imply $d(n)=[x^n]S_{n-1}+n$. Therefore
\[
T(2n,n)=[x^n]R_n=-[x^n]S_{n-1}=n-d(n),
\]
including $n=1$, where both sides are zero.
\end{proof}

The Lean module \texttt{Uchimura.lean} proves the nilpotent-ring algebra
in general, then applies it to the quotient by $X^{N+1}$.
The companion \texttt{Proof.lean} supplies the original recurrence bridge
and the factor-pair bijection with \texttt{Nat.divisors}.
Uchimura's identity and finite-product methods are established prior
mathematics; the result here is their checked connection to the two
specific triangle conjectures.

\section{A393833: palindromic proper-divisor residues}

Michel Lagneau's March 1, 2026 entry A393833~\cite{A393833} asks for the
least positive integer $k$ with exactly $2n+2$ divisors such that its
$2n+1$ proper divisors, listed in increasing order and reduced modulo
$n$, form a palindrome. The recorded conjecture is existence for every
positive $n$.

\begin{theorem}\label{thm:palindrome}
For all $m,r\ge1$ and every bound $B$, there exists $k>B$ with exactly
$r+1$ divisors whose proper-divisor residue list modulo $m$ consists of
$r$ copies of $1\bmod m$.
In particular, the minimum defining A393833 exists at every positive
index, and the eligible integers are arbitrarily large.
\end{theorem}

\begin{proof}
The classical theorem on primes congruent to one supplies a prime
$p>B$ with $p\equiv1\pmod m$. Set $k=p^r$. Its divisors, in increasing
order, are exactly $1,p,p^2,\ldots,p^r$. Thus it has $r+1$ divisors,
of which $r$ are proper. Each proper divisor is $1$ modulo $m$, and
$p^r\ge p>B$.

For the original problem take $m=n$ and $r=2n+1$. The residue list is
constant and hence a palindrome, with exactly the prescribed length and
divisor count. At $n=1$ the common residue is zero, so this boundary case
is included. Well-ordering gives a least eligible integer, while the
arbitrary bound gives infinitely many eligible integers.
\end{proof}

The source's definition and executable predicates permit constant
palindromes. The construction need not be minimal: at $n=1$ it can
produce $8$, whereas the displayed minimum is $6$. No new minimum values
are claimed. A stale example for $n=9$ in the source differs from its
corrected data; neither value is used in the proof.

Lean defines the list by sorting \texttt{Nat.properDivisors k} and
mapping reduction modulo $n$. The exact list equality, total divisor
count, proper-list length, palindrome, arbitrary-bound result, and
existence of the least witness are all proved. The prime theorem is
mathlib's existing \texttt{Nat.exists\_prime\_gt\_modEq\_one}, whose
cyclotomic proof is credited to its authors~\cite{PrimesOne}.

\section{Verification, source status, and contribution boundaries}

Each of the five packages contains a complete written proof, exact source
snapshots, a prior-work record, pinned build configuration, a build log,
and an independent audit. The six proof modules were independently
compiled in fresh directories. In the two-module package, the second
module used the freshly compiled first module.

All packages use Lean 4.19.0 and mathlib commit
\texttt{c44e0c8ee63ca166450922a373c7409c5d26b00b}.
The public theorem audits contain only the standard axioms
\texttt{propext}, \texttt{Classical.choice}, and \texttt{Quot.sound}.
There are no placeholders, additional axioms, external proof oracles,
or \texttt{native\_decide} calls. The theorem statements have their full
specified domains; finite numerical checks are not proof premises.

\begin{center}
\begin{tabular}{@{}l>{\raggedright\arraybackslash}p{4.7in}@{}}
\toprule
Package & Main checked declarations\\
\midrule
A016789 & \texttt{conjecture}, \texttt{original\_conjecture}; signed-base existence and uniqueness.\\
A300984 & \texttt{conjecture}, \texttt{strengthened\_conjecture}, \texttt{member\_ge\_676}.\\
A324017 & \texttt{fixed\_point}, \texttt{quotient\_conjecture}, \texttt{third\_claim\_false}, \texttt{corrected\_compatibility}.\\
A389524 & \texttt{row\_sum}, \texttt{diagonal}; \texttt{finite\_uchimura\_coeff} and original recurrence bridge.\\
A393833 & \texttt{conjecture}, \texttt{arbitrarily\_large\_witnesses}, \texttt{exists\_least}.\\
\bottomrule
\end{tabular}
\end{center}

The preserved primary snapshots still label the selected assertions
conjectural: A016789 revision 305; A300984 revision 15; A324017 revision
91; A389524 revision 19; and A393833 revision 30. Source checks were made
on November 26, 2026, America/Chicago. An unchanged label alone is not
evidence of novelty. Exact-statement and semantic searches, inspection
of linked literature, and checks of public proof collections were used
to distinguish related established results from the targets. Another
candidate, A049803, was excluded after an earlier published proof was
located.

The count is eight resolved assertions in five entries: seven true and
one false. The companion XOR identities, strengthened squarefree
conditions, minimum lower bound, corrected tetration relation, and
arbitrary-bound construction are not counted again. A300984's separate
frequency conjecture and A389524's weighted orphan-count conjecture
remain outside this manuscript. The work does not assert OEIS editorial
acceptance or exhaustive publication priority.

To reproduce a package with the pinned toolchain selected, run
\texttt{lake exe cache get} followed by \texttt{lake build} in its
directory. Per-package \texttt{SHA256SUMS.txt} files identify the exact
delivered artifacts. The source, audits, and this manuscript accompany
the repository publication.

\begin{thebibliography}{99}
\raggedright
\bibitem{A016789} OEIS Foundation, A016789, including G. Breysens's
July 21, 2017 conjecture; and A005351, A005352.
\url{https://oeis.org/A016789};
\url{https://oeis.org/A005351}; \url{https://oeis.org/A005352}.
\bibitem{Arndt} J. Arndt, \emph{Matters Computational}, section 1.22,
pp. 58--60. \url{https://www.jjj.de/fxt/fxtbook.pdf}.
\bibitem{A300984} M. Lagneau, OEIS A300984, March 17, 2018.
\url{https://oeis.org/A300984}.
\bibitem{A048250} OEIS A048250, including V. Jovovic's April 19, 2001
product formula. \url{https://oeis.org/A048250}.
\bibitem{A065302} OEIS A065302, including W. Kehowski's March 21, 2007
squarefree-divisor-sum argument. \url{https://oeis.org/A065302}.
\bibitem{A324017} D. Smith, OEIS A324017, March 28, 2019.
\url{https://oeis.org/A324017}.
\bibitem{Crux1981} C. W. Trigg, Problem 559 and its discussion,
\emph{Crux Mathematicorum} \textbf{7} (1981), p. 192.
\url{https://cms.math.ca/wp-content/uploads/crux-pdfs/Crux_v7n06_Jun.pdf}.
\bibitem{Crux1987} D. B. Shapiro, solution to Problem 559,
\emph{Crux Mathematicorum} \textbf{13} (1987), 291--294.
\url{https://cms.math.ca/wp-content/uploads/crux-pdfs/Crux_v13n09_Nov.pdf}.
\bibitem{Shapiro} D. B. Shapiro and S. David Shapiro,
Iterated Exponents in Number Theory, \emph{Integers} \textbf{7} (2007), A23.
Author-hosted version:
\url{https://bpb-us-w2.wpmucdn.com/u.osu.edu/dist/c/5162/files/2018/02/IE-24ed9ge.pdf}.
\bibitem{Yebra} J. Jim\'enez Urroz and J. Luis A. Yebra,
On the Equation $a^x\equiv x\pmod{b^n}$,
\emph{Journal of Integer Sequences} \textbf{12} (2009), Article 09.8.8.
\url{https://cs.uwaterloo.ca/journals/JIS/VOL12/Yebra/yebra4.pdf}.
\bibitem{A389524} M. Granvik, OEIS A389524, October 7, 2025;
target sequences A000005 and A049820.
\url{https://oeis.org/A389524}; \url{https://oeis.org/A049820}.
\bibitem{Uchimura} K. Uchimura, An identity for the divisor generating
function arising from sorting theory,
\emph{Journal of Combinatorial Theory, Series A} \textbf{31} (1981), 131--135.
\url{https://doi.org/10.1016/0097-3165(81)90009-1}.
\bibitem{Bressoud} D. M. Bressoud and M. V. Subbarao,
On Uchimura's connection between partitions and the number of divisors,
\emph{Canadian Mathematical Bulletin} \textbf{27} (1984), 143--145.
\url{https://doi.org/10.4153/CMB-1984-022-5}.
\bibitem{A393833} M. Lagneau, OEIS A393833, March 1, 2026.
\url{https://oeis.org/A393833}.
\bibitem{PrimesOne} R. Brasca and mathlib contributors,
\emph{Primes congruent to one}, pinned mathlib source.
\url{https://github.com/leanprover-community/mathlib4/blob/c44e0c8ee63ca166450922a373c7409c5d26b00b/Mathlib/NumberTheory/PrimesCongruentOne.lean}.
\end{thebibliography}
\end{document}
